%jbb, 4.7
On donne à la figure suivante le diagramme binaire solide-liquide \ce{CaCl2}/\ce{NaCl}. À l'état solide, on constate une miscibilité partielle des solides.

\begin{center}
\begin{minipage}{0.48\textwidth}
\includegraphics[width=\textwidth]{/home/jeff/Documents/IUTOrsay/1-Cours/4p13-Alliages/Cours-TD/Exos/Binaires-Lecture/MiscibilitePartielle-Ex02/figures/Fig01}
\end{minipage}\hfill\begin{minipage}{0.48\textwidth}
\captionof{figure}{Diagramme binaire isobare, à une pression de \SI{1}{\bar}, du mélange \ce{CaCl2}-\ce{NaCl}}
\end{minipage}
\end{center}

\Q Préciser la nature et la formule du composé \textbf{C}.

\Q Préciser la nature des phases dans chacun des domaines du diagramme.

\Q On considère un système renfermant \SI{0.2}{\mole} de \ce{CaCl2} à \SI{560}{\celsius}. On ajoute progressivement \ce{NaCl}. Déterminer la quantité minimale de \ce{NaCl} à ajouter pour obtenir :
\begin{itemize}
\item un mélange complètement liquide ;
\item un mélange complètement solide (préciser la nature du solide).
\end{itemize}

\paragraph*{Données}
\begin{align*}
M(\ce{Na}) &= 23 ~ \si{\gram\per\mole} ; 
M(\ce{Ca}) &=  40 ~ \si{\gram\per\mole} ; 
M(\ce{Cl}) &=  35.5  ~ \si{\gram\per\mole}.
\end{align*}
